\anonsection{ЗАКЛЮЧЕНИЕ}
В рамках выпускной квалификационной работы поставлена и решена задача
разработки компилятора базисного подмножества Рефала-5 с представлением
объектных выражений на основе раскрученной двусторонней очереди, а также
подтверждения корректности и работоспособности реализации автоматическими
тестами и сравнительным прогоном с двумя альтернативными реализациями
языка.

В исследовательской части рассмотрены базовые понятия Рефала-5 и
существующие представления объектных выражений (классическое плоское
двусвязное, компромиссное с подвешенными скобками, массивное,
векторно-списковое и дерево конкатенаций), обсуждены их компромиссы по
ключевым операциям и обоснован выбор раскрученной двусторонней очереди по
Окасаки~\cite{okasaki_pfds} как структуры, амортизированно покрывающей все
основные операции (отделение терма с концов, конкатенацию и многократное
использование значения переменной).

В конструкторской части зафиксированы функциональные и нефункциональные
требования к системе, декомпозиция на модули с привязкой к Java-пакетам,
интерфейс контейнера выражения и требования к механизму сопоставления, выделены
точки расширения и определены классы ошибок с единым форматом
диагностических сообщений.

В технологической части реализован компилятор с интерпретирующим исполнением,
выполняющий лексический,
синтаксический и семантический анализ исходного текста и формирующий
внутреннее представление программы. Раскрученная двусторонняя очередь
реализована через запечатанный интерфейс \prog{Expr<T>} и пару классов
\prog{ShallowDeque<T>} (циклический массив фиксированной ёмкости
\(CAPACITY=9\)) и \prog{DeepDeque<T>} (тройка \(\langle front, mid, back
\rangle\) с рекурсивным средним уровнем); конкатенация со всеми четырьмя
случаями Shallow/Deep $\times$ Shallow/Deep вынесена в служебный класс
\prog{ExprOps}. Исполняющий модуль реализован как цикл нормализации в
прикладном порядке с сопоставлением образца через перебор длины
\(e\)-переменной и сборкой результата через \prog{Expr.concat}. Реализован
минимальный набор встроенных функций (ввод-вывод, арифметика, работа со
строками и символами, функции высшего порядка), достаточный для тестовых
программ и для выхода рассахаривателя \prog{r5fw-desugar} проекта
\prog{refal-5-framework}~\cite{refal5_framework_github}; директива
\prog{\$EXTERN} поддерживается, что обеспечивает совместимость с этим
рассахаривателем.

В аналитической части построен автоматический набор из \(224\) тестов:
\(143\) теста JUnit~5~\cite{junit5_guide}, покрывающих лексер, парсер,
семантический анализ, сопоставление, подстановку, цикл нормализации и саму
раскрученную очередь, и \(81\) внешняя программа в рамках
\prog{AutotestsRunnerTest}~--- фикстуры из репозиториев
\prog{refal-5j} и \prog{refal-5-framework}~\cite{refal5_framework_github}.
Среднее покрытие по инструкциям байт-кода составляет около \(87\,\%\),
по веткам~--- около \(77\,\%\), что выше зафиксированных в \prog{pom.xml} порогов.
Дополнительно проведены два сравнительных прогона трёх реализаций Рефала-5
(\prog{refal5q-impl}, Refal-5~PZ~\cite{refal5_home} и Refal-5J): синтетический
тест Loop/Select и прогон рассахаривателя \prog{R5FW-Parser.ref} из проекта
\prog{refal-5-framework}~\cite{refal5_framework_github}. Замеры автоматизированы
скриптами \prog{bench/loop-select-bench.ps1} и \prog{bench/desugar-bench.ps1};
методика~--- один прогревочный запуск и три измеряемых, в отчёт идёт медиана.

Основные численные результаты: автоматический прогон \prog{mvn clean
verify} проходит без нарушений; контроль стиля Checkstyle и
покрытие выполняются выше установленных порогов. На тесте Loop/Select
на малых $N$ разработанная реализация проигрывает конкурентам из‑за
накладных расходов запуска JVM ($\approx 120$~мс), однако при росте входа
зависимость времени от $N$ остаётся линейной: при $N = 20\,000$ время
составляет 302~мс против 1059~мс у Refal-5~PZ и 1348~мс у Refal-5J
(разрыв 3,5 и 4,5~раза соответственно). На прогоне рассахаривателя все
три реализации выдают эквивалентный результат; по времени Refal-5~PZ
быстрее в 5,6~раза, Refal-5J~--- в 1,3~раза, что соответствует ожидаемому
преимуществу компиляции над интерпретацией на программе без квадратичного
паттерна.

Возможные направления дальнейшего развития системы:

\begin{itemize}
  \item Компиляция в JVM-байт-код по образцу Refal-5J: порождение для
        каждого правила отдельного Java-метода уменьшило бы постоянный
        множитель и сократило бы отставание от Refal-5J на программах
        без квадратичного паттерна, например на прогоне рассахаривателя.
  \item Поддержка модульной системы Рефала-5 (директива \prog{*\$FROM},
        квалифицированные имена). Текущая реализация поддерживает
        конкатенацию нескольких файлов через синтаксис \prog{+}, однако
        \prog{*\$FROM} и квалифицированные имена функций не реализованы.
  \item Расширение библиотеки встроенных функций до полного набора
        Рефал-5 (\prog{Br}/\prog{Dg}, шифр-функции, файловый ввод-вывод).
        Архитектура регистрации позволяет это сделать без изменений ядра.
\end{itemize}

\indent Полученная программная система (Maven-проект \prog{refal5q-impl} объёмом
\(\sim 3\,800\) строк основного кода и \(\sim 1\,700\) строк тестов) может
использоваться как самостоятельный компилятор базисного подмножества
Рефала-5 и как базовая платформа для дальнейших исследований в области
эффективных представлений объектных выражений и компиляции функциональных
языков переписывания термов.
